% math0 manuscript. Build with: latexmk -pdf paper.tex
\documentclass[11pt]{amsart}
\usepackage[margin=1.1in]{geometry}
\usepackage{amsmath,amssymb,amsthm,mathtools}
\usepackage{url}
\usepackage{hyperref}
\usepackage{xcolor}
\usepackage{listings}
\usepackage{newunicodechar}

\lstdefinelanguage{lean}{
  keywords={theorem,lemma,def,abbrev,example,fun,by,have,show,from,let,in,if,then,else,match,with,
    namespace,end,open,import,section,variable,instance,structure,where,calc,obtain,intro,exact,
    refine,rcases,rintro,simp,omega,decide,norm_num,linarith,nlinarith,positivity,induction,cases,
    private},
  sensitive=true, comment=[l]{--}, morecomment=[s]{/-}{-/}, string=[b]"}
\lstset{language=lean, basicstyle=\ttfamily\footnotesize, columns=fullflexible, keepspaces=true, breaklines=true,
  keywordstyle=\color{blue!60!black}\bfseries, commentstyle=\color{gray}, frame=single,
  extendedchars=true, inputencoding=utf8,
  literate={ℕ}{{$\mathbb{N}$}}1 {ℤ}{{$\mathbb{Z}$}}1 {ℚ}{{$\mathbb{Q}$}}1 {∀}{{$\forall$}}1
    {→}{{$\to$}}1 {≤}{{$\le$}}1 {∈}{{$\in$}}1 {∧}{{$\wedge$}}1 {∑}{{$\sum$}}1 {·}{{$\cdot$}}1
    {⁻¹}{{$^{-1}$}}2 {⟨}{{$\langle$}}1 {⟩}{{$\rangle$}}1}
% Unicode commonly found in Lean sources.
\newunicodechar{ℕ}{\ensuremath{\mathbb{N}}}
\newunicodechar{ℤ}{\ensuremath{\mathbb{Z}}}
\newunicodechar{ℚ}{\ensuremath{\mathbb{Q}}}
\newunicodechar{ℝ}{\ensuremath{\mathbb{R}}}
\newunicodechar{ℂ}{\ensuremath{\mathbb{C}}}
\newunicodechar{∀}{\ensuremath{\forall}}
\newunicodechar{∃}{\ensuremath{\exists}}
\newunicodechar{→}{\ensuremath{\to}}
\newunicodechar{←}{\ensuremath{\leftarrow}}
\newunicodechar{↔}{\ensuremath{\leftrightarrow}}
\newunicodechar{↦}{\ensuremath{\mapsto}}
\newunicodechar{¬}{\ensuremath{\neg}}
\newunicodechar{∧}{\ensuremath{\wedge}}
\newunicodechar{∨}{\ensuremath{\vee}}
\newunicodechar{≤}{\ensuremath{\le}}
\newunicodechar{≥}{\ensuremath{\ge}}
\newunicodechar{≠}{\ensuremath{\ne}}
\newunicodechar{∈}{\ensuremath{\in}}
\newunicodechar{∉}{\ensuremath{\notin}}
\newunicodechar{⊆}{\ensuremath{\subseteq}}
\newunicodechar{∩}{\ensuremath{\cap}}
\newunicodechar{∪}{\ensuremath{\cup}}
\newunicodechar{∑}{\ensuremath{\sum}}
\newunicodechar{∏}{\ensuremath{\prod}}
\newunicodechar{⟨}{\ensuremath{\langle}}
\newunicodechar{⟩}{\ensuremath{\rangle}}
\newunicodechar{·}{\ensuremath{\cdot}}
\newunicodechar{∘}{\ensuremath{\circ}}
\newunicodechar{⁻}{\ensuremath{^{-}}}
\newunicodechar{¹}{\ensuremath{^{1}}}
\newunicodechar{²}{\ensuremath{^{2}}}
\newunicodechar{₀}{\ensuremath{_{0}}}
\newunicodechar{₁}{\ensuremath{_{1}}}
\newunicodechar{₂}{\ensuremath{_{2}}}
\newunicodechar{α}{\ensuremath{\alpha}}
\newunicodechar{β}{\ensuremath{\beta}}
\newunicodechar{ε}{\ensuremath{\varepsilon}}
\newunicodechar{δ}{\ensuremath{\delta}}
\newunicodechar{π}{\ensuremath{\pi}}
\newunicodechar{σ}{\ensuremath{\sigma}}
\newunicodechar{φ}{\ensuremath{\varphi}}
\newunicodechar{∞}{\ensuremath{\infty}}
\newunicodechar{⌊}{\ensuremath{\lfloor}}
\newunicodechar{⌋}{\ensuremath{\rfloor}}
\newunicodechar{⌈}{\ensuremath{\lceil}}
\newunicodechar{⌉}{\ensuremath{\rceil}}
\newunicodechar{∣}{\ensuremath{\mid}}
\newunicodechar{≡}{\ensuremath{\equiv}}
\newunicodechar{×}{\ensuremath{\times}}
\newunicodechar{⧸}{\ensuremath{/}}
\newunicodechar{‖}{\ensuremath{\|}}
\newunicodechar{√}{\ensuremath{\surd}}
\newunicodechar{∅}{\ensuremath{\emptyset}}
\newunicodechar{⦃}{\ensuremath{\{\!|}}
\newunicodechar{⦄}{\ensuremath{|\!\}}}
\newunicodechar{𝓝}{\ensuremath{\mathcal{N}}}
\newunicodechar{≫}{\ensuremath{\gg}}
\newunicodechar{≪}{\ensuremath{\ll}}

\theoremstyle{plain}
\newtheorem{theorem}{Theorem}
\newtheorem{lemma}[theorem]{Lemma}
\newtheorem{proposition}[theorem]{Proposition}
\newtheorem{corollary}[theorem]{Corollary}
\theoremstyle{definition}
\newtheorem{definition}[theorem]{Definition}
\newtheorem{remark}[theorem]{Remark}

\emergencystretch=3em
\allowdisplaybreaks
% Lean identifiers in running text (breakable at dots and underscores).
\expandafter\def\expandafter\UrlBreaks\expandafter{\UrlBreaks\do\_\do\.}
\DeclareUrlCommand\decl{\def\UrlFont{\normalfont\ttfamily}}
\newcommand{\N}{\mathbb{N}}
\newcommand{\Z}{\mathbb{Z}}
\newcommand{\Q}{\mathbb{Q}}
\newcommand{\ib}[1]{[\,#1\,]}
\newcommand{\lb}{\char123}
\newcommand{\rb}{\char125}

\title[Schulte's conjectures on OEIS A281000]{Antidiagonal sums and the matrix inverse of the triangle $\binom{2n+1}{2k+1}C_{n-k}$: proof of Schulte's conjectures on OEIS A281000}
\author{math0}
\address{Produced by Claude (an AI model developed by Anthropic) in the math0 project}
\date{October 25, 2026}


\begin{document}
\begin{abstract}
Let $T(n,k)=\binom{2n+1}{2k+1}C_{n-k}$ for $0\le k\le n$, where $C_m$ are the Catalan numbers (OEIS A281000).
We prove the conjectures that Werner Schulte recorded in this entry in 2017: the antidiagonal sums of $T$ are the little Schr\"oder numbers $s(n+1)$; the inverse of the lower unitriangular matrix $T$ is $\bigl(T(n,k)\,b(n+1-k)\bigr)$, where $b$ (A103365) is the first column of the inverse Narayana matrix; and, in a corrected reading of a conjecture that as printed refers to the undefined term A103365$(0)$, $\sum_kT(n,k)(2k+1)C_k\,b(k+1)=0^n$.
The proofs are finite binomial-coefficient identities, and the math0 referees found no earlier proof, so we classify the result as new; it is routine, however, and can also be derived from generating-function facts recorded in related OEIS entries.
The proof is verified in Lean~4 against the frozen, fidelity-reviewed statement that math0 formalized from the OEIS entry.
\end{abstract}
\maketitle

\section{Introduction}

Let $C_m=\binom{2m}{m}/(m+1)$ be the Catalan numbers (A000108 \cite{A000108}). The triangle
\[
  T(n,k)=\binom{2n+1}{2k+1}\binom{2n-2k}{n-k}\Big/(n+1-k)=\binom{2n+1}{2k+1}C_{n-k}\qquad(0\le k\le n),
\]
with rows $1;\ 3,1;\ 10,10,1;\ 35,70,21,1;\dots$, is the entry A281000 of the OEIS \cite{A281000}, submitted by Werner Schulte on January 12, 2017. We put $T(n,k)=0$ for $k>n$, so $T$ is an infinite lower triangular matrix with unit diagonal. As the entry notes, $T(n,k)=M(2n+1,2k+1)$, where $M(n,k)=\binom nk C_{(n-k)/2}$ ($0$ if $n-k$ is odd) counts Motzkin paths of length $n$ with $k$ level steps (A097610 \cite{A097610}). Schulte recorded three conjectures in the entry, on the antidiagonal sums of $T$, a weighted row sum, and the inverse matrix (Section~\ref{sec:statement}); in revision \#17 (November 25, 2018), the version in our local mirror of the OEIS, they are still labelled ``Conjectures''. We prove all three (Theorem~\ref{thm:main}), the second in the corrected reading of Remark~\ref{rem:reading}.

\subsection*{Prior work and status}
\begin{itemize}
\item The math0 literature referees checked A281000 and the related entries A280580, A103365, A103364, A238390, A097610, A001003, A006318, A126120 and A180874 (no entry other than A281000 mentions A281000), and searched arXiv (including recent papers devoted to proofs of OEIS conjectures), MathOverflow, the OEIS wiki and the SeqFan list. They found no proof of any of the conjectures.
\item \emph{(3) and (2) from known formulas.} Tom Copeland states in A238390 \cite{A238390} (January 30, 2016) that the matrices $\binom nk A126120(n-k)$ (A126120 \cite{A126120} is the aerated Catalan sequence) and $\binom nk m(n-k)$, with $m(2j)=(-1)^jA238390(j)$ and $m(2j+1)=0$, are mutually inverse. The first is $M$; both vanish unless $n-k$ is even, so $T^{-1}$ is the odd--odd block of the second, $T^{-1}(n,k)=\binom{2n+1}{2k+1}(-1)^{n-k}A238390(n-k)$. Comparing the e.g.f.\ $x/J_1(2x)$ of A238390 with P.~D.~Hanna's g.f.\ in A103365 \cite{A103365} (2009) gives $(-1)^jA238390(j)=C_j\,b(j+1)$, whence (3), and (2) is its first column. This derivation was pointed out by the referees; we did not find it in print.
\item \emph{(1) from known formulas.} Paul Barry notes in A097610 (April 21, 2010) that the diagonal sums $\sum_kM(m-k,k)$ are the aerated large Schr\"oder numbers $r$ (A006318 \cite{A006318}). For $m=2n+2$ the terms with $k$ odd form the $n$-th antidiagonal sum of $T$; since $\sum_k(-1)^kM(m-k,k)=0$ for $m\ge1$ (a generating-function computation not in the entry) and $r(n)=2s(n)$ for $n\ge1$, this gives (1).
\end{itemize}
We therefore classify the result as new, with the caveat that it is a routine consequence of known generating-function facts. Our proof uses no generating functions: it consists of finite binomial identities, chosen because they are easy to formalize. Remark~\ref{rem:toeplitz} explains the structure behind the proof of (3); it is close to Peter Bala's description (2013) of the inverse Narayana matrix A103364 \cite{A103364} as a generalized Riordan array.

\section{Statement}\label{sec:statement}

The FORMULA section of A281000 \cite{A281000} (offset $0$, $0\le k\le n$; entry by Werner Schulte, January 12, 2017) contains, verbatim,
\begin{quote}\raggedright\footnotesize
\texttt{Conjectures:}\\
\texttt{\ \ (1) Antidiagonal sums equal A001003(n+1);}\\
\texttt{\ \ (2) Sum\_\lb k=0..n\rb\ (-1)\^{}k*T(n,k)*(2*k+1)*A000108(k)*A103365(k) = A000007(n);}\\
\texttt{\ \ (3) Matrix inverse equals T(n,k)*A103365(n+1-k).}
\end{quote}
Here A000007$(n)=0^n$. Let $N(n,j)=\binom{n-1}{j-1}\binom n{j-1}/j$ ($1\le j\le n$) be the Narayana numbers (A001263 \cite{A001263}); the matrix $(N(n,j))_{n,j\ge1}$ is lower unitriangular. Let $s(n)$ be the little Schr\"oder numbers $1,1,3,11,45,\dots$ (A001003 \cite{A001003}); we use B.~Cloitre's formula from that entry (2003, reindexed), $s(n)=\sum_{j=1}^{n}2^{j-1}N(n,j)$ for $n\ge1$, as the definition. Let $b(m)$, $m\ge1$, be A103365 \cite{A103365}, ``the first column of triangle A103364, which equals the matrix inverse of the Narayana triangle (A001263)'', i.e.\ the solution of $\sum_{j=1}^{n}N(n,j)\,b(j)=\ib{n=1}$ $(n\ge1)$: $1,-1,2,-7,39,-321,\dots$

\begin{theorem}\label{thm:main}
\begin{enumerate}
\item[(1)] For every $n\ge0$,  $\displaystyle\sum_{k=0}^{\lfloor n/2\rfloor}T(n-k,k)=s(n+1)$;
\item[(2)] for every $n\ge0$, $\displaystyle\sum_{k=0}^{n}T(n,k)\,(2k+1)\,C_k\,b(k+1)=0^n$;
\item[(3)] for every $N\ge0$, the matrix $(T(i,j))_{0\le i,j<N}$ over $\Q$ is invertible, with inverse $(U(i,j))_{0\le i,j<N}$, where $U(i,j)=T(i,j)\,b(i+1-j)$ for $j\le i$ and $U(i,j)=0$ for $j>i$.
\end{enumerate}
\end{theorem}
As $T$ and $U$ are lower triangular, (3) is equivalent to $T^{-1}=U$ for the infinite matrices.

\begin{remark}[The reading of (2)]\label{rem:reading}
As printed, (2) is not well defined: at $k=0$ it uses A103365$(0)$, but A103365 has offset $1$. With the index shift A103365$(k)\mapsto b(k+1)$ alone it is false at $n=1$, where the sum is $3+3=6$, because $(-1)^k$ cancels the alternating signs of $b$; giving A103365$(0)$ the value $0$ or $1$ fails as well. We read $(-1)^k$A103365$(k)$ as $b(k+1)$, i.e.\ A103365 as the unsigned sequence indexed from $0$ (for the listed terms, $b(k+1)$ has sign $(-1)^k$). This is the convention of Hanna's g.f.\ in A103365, where $(-1)^na(n)$ stands for $b(n+1)$, and of Schulte's analogous conjecture $\sum_k(-1)^k\binom{2n}{2k}C_{n-k}A238390(k)=0^n$ in A280580 \cite{A280580} (January 5, 2017), with A238390 unsigned of offset $0$. So read, (2) is the column $k=0$ of $T\cdot U=I$, since $T(k,0)=(2k+1)C_k$; it does not change the strength of the statement, which is equivalent to (1) and (3).
\end{remark}

\subsection*{Lean formalization}
The statement was formalized by math0 in \path{lean/Math0/Statements/OEIS/A281000.lean}, reviewed for fidelity, and frozen (SHA-256 \texttt{4230cf53\ldots c030}) before any proof attempt. Its definitions and the statement read (verbatim; we omit the import, comments, docstrings, all theorems, and the definition \decl{a} of the triangle read by rows, used only in the data check):
\begin{lstlisting}
namespace Math0.Statements.A281000
def T (n k : ℕ) : ℕ :=
  Nat.choose (2 * n + 1) (2 * k + 1) * Nat.choose (2 * n - 2 * k) (n - k) / (n + 1 - k)
def narayana (n k : ℕ) : ℕ :=
  Nat.choose (n - 1) (k - 1) * Nat.choose n (k - 1) / k
def A001003 (n : ℕ) : ℕ :=
  if n = 0 then 1 else ∑ j ∈ Finset.Icc 1 n, 2 ^ (j - 1) * narayana n j
def invNarayanaCol : ℕ → List ℤ
  | 0 => []
  | m + 1 =>
    let l := invNarayanaCol m
    l ++ [(if m = 0 then 1 else 0) -
      ((List.range m).map fun j => (narayana (m + 1) (j + 1) : ℤ) * l.getD j 0).sum]
def A103365 (n : ℕ) : ℤ :=
  (invNarayanaCol n).getD (n - 1) 0
def conj1 : Prop :=
  ∀ n : ℕ, ∑ k ∈ Finset.range (n / 2 + 1), T (n - k) k = A001003 (n + 1)
def conj2 : Prop :=
  ∀ n : ℕ, ∑ k ∈ Finset.range (n + 1),
    (T n k : ℤ) * (2 * k + 1) * (catalan k : ℤ) * A103365 (k + 1) = (0 : ℤ) ^ n
def conj3 : Prop :=
  ∀ N : ℕ, (Matrix.of fun i j : Fin N => (T i j : ℚ))⁻¹ =
    Matrix.of fun i j : Fin N => (T i j : ℚ) * (A103365 ((i : ℕ) + 1 - (j : ℕ)) : ℚ)
def statement : Prop :=
  conj1 ∧ conj2 ∧ conj3
end Math0.Statements.A281000
\end{lstlisting}
\emph{Encoding.} \texttt{T} is the formula in the name of the entry, computed in $\N$ with truncated subtraction and floor division; it equals $\binom{2n+1}{2k+1}C_{n-k}$ for $k\le n$ and $0$ for $k>n$ (Lemma~\ref{lem:closed}). Likewise \texttt{narayana} is $N$, with exact division. \texttt{A001003} is Cloitre's formula. \texttt{invNarayanaCol m} is the list $[b(1),\dots,b(m)]$ computed by forward substitution, $b(m+1)=\ib{m=0}-\sum_{j<m}N(m+1,j+1)\,b(j+1)$ (row $m+1$ of the defining system, as $N(m+1,m+1)=1$), so \texttt{A103365 n} $=b(n)$ for $n\ge1$; \texttt{A103365 0} $=0$ is a junk value. In \texttt{conj3} the truncated $i+1-j$ only occurs multiplied by \texttt{T i j} $=0$ when $j>i$. Mathlib's \texttt{catalan} is A000108, and \texttt{⁻¹} is Mathlib's matrix inverse (which would be $0$ for a singular matrix; here $\det=1$).

\emph{Data checks.} By \texttt{decide +kernel}, the statement file checks that the triangle read by rows reproduces all $55$ listed terms of A281000 (\decl{a_data}), that \texttt{A001003} reproduces all $26$ listed terms $s(0),\dots,s(25)$ (\decl{A001003_data}), and that \texttt{A103365} reproduces all $20$ listed terms $b(1),\dots,b(20)$ (\decl{A103365_data}). It also checks \texttt{conj1} for $n\le24$, \texttt{conj2} for $n\le15$, and the entries of $T\cdot U$ for indices $\le14$; these are transcription checks only.

\section{Proof}\label{sec:proof}

The proof follows \path{lean/Math0/Frontier/OEIS/A281000.lean}; Lean names are given in parentheses. We write $\ib P$ for $1$ if $P$ holds and $0$ otherwise, and $\binom ab=0$ for $b>a$.

\subsection{Closed forms}

\begin{lemma}\label{lem:closed}
For all $n,k,p,M,i,j\ge0$:
\begin{enumerate}
\item[(a)] $T(n,k)=0$ if $n<k$ \quad(\decl{T_eq_zero});
\item[(b)] $T(k+p,k)=\binom{2k+2p+1}{2k+1}C_p$, in particular $T(k,0)=(2k+1)C_k$ \quad(\decl{T_add}, \decl{T_zero});
\item[(c)] $T(k+p,k)\,(2k+1)!\,p!\,(p+1)!=(2k+2p+1)!$ \quad(\decl{T_mul_fact});
\item[(d)] $(i+1)\mid\binom Mi\binom{M+1}i$ and $N(M+1,i+1)\,(i+1)=\binom Mi\binom{M+1}i$; hence $N(M+1,M+1)=1$ \quad(\decl{nar_dvd}, \decl{nar_mul});
\item[(e)] $N(i+j+1,j+1)\,(j+1)!\,j!\,i!\,(i+1)!=(i+j+1)!\,(i+j)!$ \quad(\decl{nar_mul_fact}).
\end{enumerate}
\end{lemma}

\begin{proof}
(a) $\binom{2n+1}{2k+1}=0$, so the numerator vanishes. (b) For $n=k+p$ the numerator is $\binom{2k+2p+1}{2k+1}\binom{2p}{p}$ and $\binom{2p}p=(p+1)C_p$, so the division by $n+1-k=p+1$ is exact. (c) $C_p\,p!\,(p+1)!=\binom{2p}{p}p!\,p!=(2p)!$, and $\binom{2k+2p+1}{2k+1}(2k+1)!\,(2p)!=(2k+2p+1)!$. (d) Let $x=\binom Mi\binom{M+1}i$. From $(M+1)\binom Mi=(i+1)\binom{M+1}{i+1}$ and $(M+2)\binom{M+1}i=(i+1)\binom{M+2}{i+1}$, $i+1$ divides $(M+1)x$ and $(M+2)x$, hence their difference $x$; and $N(M+1,i+1)=x/(i+1)$ by definition. For $i=M$, $x=M+1$. (e) Take $M=i+j$ and index $j$ in (d), and multiply by $j!\,i!\cdot j!\,(i+1)!$, using $\binom{i+j}{j}j!\,i!=(i+j)!$ and $\binom{i+j+1}{j}j!\,(i+1)!=(i+j+1)!$.
\end{proof}

\begin{lemma}[\decl{nar_sum}]\label{lem:b}
For every $m\ge0$, $\sum_{j=0}^{m}N(m+1,j+1)\,b(j+1)=\ib{m=0}$.
\end{lemma}

\begin{proof}
By induction, \texttt{invNarayanaCol m} has length $m$ and is a prefix of \texttt{invNarayanaCol (m+1)}, so its entry $j<m$ is $b(j+1)$ (\decl{invNarayanaCol_length}, \decl{invNarayanaCol_getD}). The last entry of \texttt{invNarayanaCol (m+1)} is $b(m+1)$, so $b(m+1)=\ib{m=0}-\sum_{j<m}N(m+1,j+1)\,b(j+1)$ (\decl{A103365_succ}). Add $N(m+1,m+1)\,b(m+1)=b(m+1)$ (Lemma~\ref{lem:closed}(d)).
\end{proof}

\subsection{The inverse: conjectures (3) and (2)}

\begin{lemma}[\decl{term_eq}]\label{lem:factor}
For $k\ge0$ and $0\le j\le p$, in $\Q$,
\[
  T(k+p,k+j)\,T(k+j,k)=F(k,p)\,N(p+1,j+1),\qquad F(k,p)=\frac{(2k+2p+1)!}{(2k+1)!\,(p+1)!\,p!}.
\]
\end{lemma}

\begin{proof}
Write $p=j+i$. By Lemma~\ref{lem:closed}(c), applied to $(k+j,i)$ and $(k,j)$, the product on the left is
\[
  \frac{(2k+2p+1)!}{(2k+2j+1)!\,i!\,(i+1)!}\cdot\frac{(2k+2j+1)!}{(2k+1)!\,j!\,(j+1)!}=\frac{(2k+2p+1)!}{(2k+1)!\,i!\,(i+1)!\,j!\,(j+1)!},
\]
and by Lemma~\ref{lem:closed}(e), $F(k,p)\,N(p+1,j+1)=F(k,p)\,(p+1)!\,p!/\bigl((j+1)!\,j!\,i!\,(i+1)!\bigr)$ is the same.
\end{proof}

\begin{lemma}[\decl{main}, \decl{main'}]\label{lem:TU}
For all $n,k\ge0$ and $L>n$, in $\Q$, with $l+1-k$ truncated at $0$ and $b(0)=0$ (as in Lean),
\[
  \sum_{l=0}^{L-1}T(n,l)\,T(l,k)\,b(l+1-k)=\ib{n=k}.
\]
\end{lemma}

\begin{proof}
If $k>n$, every term vanishes by Lemma~\ref{lem:closed}(a): $T(l,k)=0$ for $l\le n$ and $T(n,l)=0$ for $l>n$. Let $n=k+p$. The terms with $l>n$ or $l<k$ vanish for the same reason. Writing $l=k+j$ with $0\le j\le p$, so that $l+1-k=j+1$, Lemmas~\ref{lem:factor} and~\ref{lem:b} turn the sum into $F(k,p)\sum_{j=0}^{p}N(p+1,j+1)\,b(j+1)=F(k,p)\,\ib{p=0}=\ib{p=0}$, as $F(k,0)=1$.
\end{proof}

\begin{proof}[Proof of Theorem~\ref{thm:main}(3) \textup{(\decl{conj3_proof})}]
Let $T_N=(T(i,j))_{i,j<N}$ and $U_N=(T(i,j)\,b(i+1-j))_{i,j<N}$. The $(i,j)$ entry of $T_NU_N$ is the sum of Lemma~\ref{lem:TU} with $n=i<L=N$ and $k=j$, so $T_NU_N=I_N$. A right inverse of a square matrix over a commutative ring is its inverse (Mathlib's \decl{Matrix.inv_eq_right_inv}), so $T_N^{-1}=U_N$.
\end{proof}

\begin{proof}[Proof of Theorem~\ref{thm:main}(2) \textup{(\decl{conj2_proof})}]
By Lemma~\ref{lem:closed}(b), $T(k,0)\,b(k+1-0)=(2k+1)\,C_k\,b(k+1)$, so the sum in (2) is that of Lemma~\ref{lem:TU} with $k=0$, $L=n+1$, and equals $\ib{n=0}=0^n$ in $\Q$, hence in $\Z$.
\end{proof}

\begin{remark}\label{rem:toeplitz}
Let $c_m=1/(m!\,(m+1)!)$. Lemma~\ref{lem:closed}(c),(e) say that $T(n,k)=\frac{(2n+1)!}{(2k+1)!}c_{n-k}$ and $N(n,j)=\frac{n!\,(n-1)!}{j!\,(j-1)!}c_{n-j}$: both matrices are diagonal conjugates of the Toeplitz matrix $(c_{n-k})$. Their inverses are the same conjugates of $(e_{n-k})$, where $\sum_me_mz^m=(\sum_mc_mz^m)^{-1}$. Comparing first columns gives $b(m+1)=(m+1)!\,m!\,e_m$, hence $T^{-1}(n,k)=\frac{(2n+1)!}{(2k+1)!}e_{n-k}=T(n,k)\,b(n+1-k)$. Lemma~\ref{lem:TU} is a finite, entrywise form of this argument. As $\sum_mc_mz^m=I_1(2\sqrt z)/\sqrt z$, this is also the source of the Bessel-function formulas in A103365 and A238390.
\end{remark}

\subsection{The antidiagonal sums: conjecture (1)}

\begin{lemma}\label{lem:binom}
For all $i,K,k,p\ge0$:
\begin{enumerate}
\item[(a)] if $i<2K$, then $\sum_{k=0}^{K-1}\binom{i+1}{2k+1}=2^i$ \quad(\decl{odd_sum});
\item[(b)] $\sum_{i=0}^{2k+p}\binom{2k+p}{i}\binom{2k+p+1}{i}\binom{i}{2k}=\binom{2k+p}{2k}\binom{2k+2p+1}{p+1}$ \quad(\decl{vandermonde_step});
\item[(c)] $\binom{2k+p}{2k}\binom{2k+2p+1}{p+1}=(2k+1)\,T(k+p,k)$ \quad(\decl{claimC}).
\end{enumerate}
\end{lemma}

\begin{proof}
(a) By Pascal's rule $\binom{i+1}{2k+1}=\binom i{2k}+\binom i{2k+1}$, the sum is $\sum_{t=0}^{2K-1}\binom it=\sum_{t=0}^{i}\binom it=2^i$, as $2K-1\ge i$. (b) Terms with $i<2k$ vanish. For $i=2k+t$, $0\le t\le p$: $\binom{2k+p}{2k+t}\binom{2k+t}{2k}=\binom{2k+p}{2k}\binom pt$ (Mathlib's \decl{Nat.choose_mul}) and $\binom{2k+p+1}{2k+t}=\binom{2k+p+1}{p+1-t}$; by Vandermonde's identity (\decl{Nat.add_choose_eq}; the term $t=p+1$ is $0$), $\sum_{t=0}^{p}\binom pt\binom{2k+p+1}{p+1-t}=\binom{2k+2p+1}{p+1}$. (c) Multiply by $(2k)!\,p!\,(p+1)!$: the left side becomes $\binom{2k+2p+1}{p+1}(p+1)!\,(2k+p)!=(2k+2p+1)!$, and so does the right side, by Lemma~\ref{lem:closed}(c).
\end{proof}

\begin{lemma}[\decl{claimA}]\label{lem:A}
For all $k,p\ge0$, with $M=2k+p$, $T(k+p,k)=\sum_{i=0}^{M}N(M+1,i+1)\binom{i+1}{2k+1}$.
\end{lemma}

\begin{proof}
By Lemma~\ref{lem:closed}(d) and $(i+1)\binom i{2k}=(2k+1)\binom{i+1}{2k+1}$, after cancelling $i+1$,
$(2k+1)\,N(M+1,i+1)\binom{i+1}{2k+1}=\binom Mi\binom{M+1}i\binom i{2k}$. Sum over $0\le i\le M$, apply Lemma~\ref{lem:binom}(b),(c), and cancel $2k+1$.
\end{proof}

\begin{proof}[Proof of Theorem~\ref{thm:main}(1) \textup{(\decl{conj1_proof})}]
For $0\le k\le\lfloor n/2\rfloor$, Lemma~\ref{lem:A} with $p=n-2k$ (so $k+p=n-k$, $M=n$) gives $T(n-k,k)=\sum_{i=0}^{n}N(n+1,i+1)\binom{i+1}{2k+1}$. Summing over $k$, exchanging the sums, and applying Lemma~\ref{lem:binom}(a) with $K=\lfloor n/2\rfloor+1>n/2\ge i/2$,
\[
  \sum_{k=0}^{\lfloor n/2\rfloor}T(n-k,k)=\sum_{i=0}^{n}N(n+1,i+1)\sum_{k=0}^{\lfloor n/2\rfloor}\binom{i+1}{2k+1}=\sum_{i=0}^{n}2^i\,N(n+1,i+1)=s(n+1).\qedhere
\]
\end{proof}

\section{Machine verification}\label{sec:lean}

The proof is the file \path{lean/Math0/Frontier/OEIS/A281000.lean} in the math0 repository, \url{https://github.com/Foundation-0/math0}. It imports the frozen statement file and the audit commands in \path{lean/Math0/Audit.lean}. The main declaration is \decl{Math0.Frontier.A281000.result}, of type \decl{statement}, the conjunction of \decl{conj1_proof}, \decl{conj2_proof} and \decl{conj3_proof}; all auxiliary declarations are \texttt{private}, and the main ones are named in Section~\ref{sec:proof}. Neither file uses \texttt{sorry}, \texttt{native\_decide}, custom axioms or \texttt{set\_option}; the proof file uses no \texttt{decide}, and the kernel evaluations of the statement file's data checks are not used by the proof. The file ends with
\begin{lstlisting}
#math0_settles Math0.Frontier.A281000.result Math0.Statements.A281000.statement
\end{lstlisting}
This command (in \path{lean/Math0/Audit.lean}, for statements that math0 formalizes itself) fails unless \decl{statement} is a \texttt{Prop}-valued definition without universe parameters; \decl{statement} and \decl{result} depend on no axioms besides \decl{propext}, \decl{Classical.choice} and \decl{Quot.sound} (rejecting \decl{sorryAx}, \decl{Lean.ofReduceBool} from \texttt{native\_decide}, and custom axioms); and the type of \decl{result} is syntactically \decl{statement} or $\neg$\,\decl{statement}, up to metadata, with no unfolding. It reports
\begin{lstlisting}[basicstyle=\ttfamily\scriptsize]
math0 audit: 'Math0.Statements.A281000.statement' depends on axioms [Classical.choice, Quot.sound, propext]
math0 audit: 'Math0.Frontier.A281000.result' depends on axioms [Classical.choice, Quot.sound, propext]
math0 statement check: 'Math0.Frontier.A281000.result' proves 'Math0.Statements.A281000.statement'
\end{lstlisting}
We use Lean~4 \cite{Lean4} \texttt{v4.33.1} and Mathlib \cite{Mathlib} \texttt{v4.33.1}, pinned in \path{lean/}; the repository also pins formal-conjectures \cite{FC} at commit \texttt{c77fdd4}, which this result does not use, as its statement was formalized by math0 from the OEIS entry. To reproduce, run \path{scripts/verify.sh} from the repository root: it fetches the Mathlib cache, builds all math0 results with \texttt{lake build} (re-running every audit command) and checks that the pinned \texttt{formal\_conjectures} and \texttt{mathlib} checkouts are unmodified. The single file can be checked with \texttt{lake env lean Math0/Frontier/OEIS/A281000.lean} in \path{lean/}; on our machine this took about 20 seconds, including loading Mathlib.

{\footnotesize
\begin{thebibliography}{99}

\bibitem{FC}
M.~Firsching, P.~Lezeau, S.~Mercuri, M.~Z. Horv\'ath, Y.~Dillies, C.~S\"onne, E.~Wieser, F.~Zhang, T.~Hubert, B.~Ag\"uera y Arcas, and P.~Kohli, \emph{Formal Conjectures: An Open and Evolving Benchmark for Verified Discovery in Mathematics}, arXiv:2605.13171, 2026. Repository: \url{https://github.com/google-deepmind/formal-conjectures}.

\bibitem{Mathlib}
The mathlib Community, \emph{The Lean mathematical library}, in: Proceedings of the 9th ACM SIGPLAN International Conference on Certified Programs and Proofs (CPP 2020), ACM, 2020, 367--381.

\bibitem{Lean4}
L.~de Moura and S.~Ullrich, \emph{The Lean 4 theorem prover and programming language}, in: Automated Deduction -- CADE 28, Lecture Notes in Comput. Sci. \textbf{12699}, Springer, 2021, 625--635.

\bibitem{A000108}
The OEIS Foundation Inc., \emph{Catalan numbers: $C(n)=\binom{2n}{n}/(n+1)=(2n)!/(n!(n+1)!)$}, entry A000108 (N.~J.~A. Sloane) in The On-Line Encyclopedia of Integer Sequences, \url{https://oeis.org/A000108}.

\bibitem{A001003}
The OEIS Foundation Inc., \emph{Schroeder's second problem (generalized parentheses); also called super-Catalan numbers or little Schroeder numbers}, entry A001003 (N.~J.~A. Sloane) in The On-Line Encyclopedia of Integer Sequences, \url{https://oeis.org/A001003}.

\bibitem{A001263}
The OEIS Foundation Inc., \emph{Triangle of Narayana numbers $T(n,k)=C(n-1,k-1)C(n,k-1)/k$ with $1\le k\le n$, read by rows}, entry A001263 (N.~J.~A. Sloane) in The On-Line Encyclopedia of Integer Sequences, \url{https://oeis.org/A001263}.

\bibitem{A006318}
The OEIS Foundation Inc., \emph{Large Schr\"oder numbers}, entry A006318 (N.~J.~A. Sloane) in The On-Line Encyclopedia of Integer Sequences, \url{https://oeis.org/A006318}.

\bibitem{A097610}
The OEIS Foundation Inc., \emph{Triangle read by rows: $T(n,k)$ is number of Motzkin paths of length $n$ and having $k$ horizontal steps}, entry A097610 (E.~Deutsch, 2004) in The On-Line Encyclopedia of Integer Sequences, \url{https://oeis.org/A097610}.

\bibitem{A103364}
The OEIS Foundation Inc., \emph{Matrix inverse of the Narayana triangle A001263}, entry A103364 (P.~D. Hanna, 2005) in The On-Line Encyclopedia of Integer Sequences, \url{https://oeis.org/A103364}.

\bibitem{A103365}
The OEIS Foundation Inc., \emph{First column of triangle A103364, which equals the matrix inverse of the Narayana triangle (A001263)}, entry A103365 (P.~D. Hanna, 2005) in The On-Line Encyclopedia of Integer Sequences, \url{https://oeis.org/A103365}.

\bibitem{A126120}
The OEIS Foundation Inc., \emph{Catalan numbers (A000108) interpolated with 0's}, entry A126120 (Ph.~Del\'eham, 2007) in The On-Line Encyclopedia of Integer Sequences, \url{https://oeis.org/A126120}.

\bibitem{A238390}
The OEIS Foundation Inc., \emph{E.g.f.: $x/\mathrm{BesselJ}(1,2x)$ (even powers only)}, entry A238390 (V.~Kotesovec, 2014) in The On-Line Encyclopedia of Integer Sequences, \url{https://oeis.org/A238390}.

\bibitem{A280580}
The OEIS Foundation Inc., \emph{Triangle read by rows: $T(n,k)=\binom{2n}{2k}\binom{2n-2k}{n-k}/(n+1-k)$ for $0\le k\le n$}, entry A280580 (W.~Schulte, 2017) in The On-Line Encyclopedia of Integer Sequences, \url{https://oeis.org/A280580}.

\bibitem{A281000}
The OEIS Foundation Inc., \emph{Triangle read by rows: $T(n,k)=\binom{2n+1}{2k+1}\binom{2n-2k}{n-k}/(n+1-k)$ for $0\le k\le n$}, entry A281000 (W.~Schulte, 2017) in The On-Line Encyclopedia of Integer Sequences, \url{https://oeis.org/A281000}; revision \#17, November 25, 2018.

\end{thebibliography}
}
\end{document}
